\documentclass[11pt]{article}
\usepackage{mathblatt}
% Probe Serie 08.10.2026: Übersichtsblatt Kurvenuntersuchung, Fassung 3.
% Jeder Eintrag ist ein Blatt; Bereiche sind Überschriften; Anker nur, wo der Name nicht reicht.
\newcommand{\kennung}{W7Q}
\newcommand{\bereich}[1]{\par\addvspace{14pt}\noindent{\large\bfseries #1}\par\nobreak\vspace{4pt}}
\newcommand{\bl}[2]{\par\noindent\makebox[1.6em][l]{$\square$}#1\ifblank{#2}{}{\hspace{0.6em}{\small\color{mbgrau}#2}}\par\vspace{3pt}}
\newcommand{\rand}{\hspace{0.6em}{\small\color{mbgrau}Rand}}
\begin{document}
\pfheftstil
\fancyfoot[C]{{\footnotesize\color{mbgrau}\kennung\hfill\thepage}}
\pfheftkopf{Kurvenuntersuchung}{Abi GK \textperiodcentered{} Übersicht}

\bereich{Eigenschaften am Term}
\bl{Verschieben und spiegeln}{$f(x-2)+3$}
\bl{Für große $x$ \textperiodcentered{} Grenzverhalten}{$\lim\limits_{x\to\infty} f(x)$}
\bl{Symmetrie\rand}{}
\bl{Nullstellen}{$f(x)=0$}

\bereich{Ableitung}
\bl{Steigen und fallen \textperiodcentered{} Monotonie}{}
\bl{Hoch- und Tiefpunkt \textperiodcentered{} Extrempunkte}{}
\bl{Wendepunkt \textperiodcentered{} Krümmung}{}
\bl{Graph skizzieren\rand}{}

\bereich{Anwenden}
\bl{Maße am Graphen}{Höhe, Breite, Maßstab}
\bl{Größte Fläche \textperiodcentered{} Extremalproblem\rand}{}
\bl{Schnittpunkte von zwei Graphen}{}
\bl{Funktionsgleichung aufstellen\rand}{}
\end{document}
